% Created 2026-03-16 Mon 17:29
% Intended LaTeX compiler: lualatex
\documentclass[11pt]{article}
\usepackage{amsmath}
\usepackage{fontspec}
\usepackage{graphicx}
\usepackage{longtable}
\usepackage{wrapfig}
\usepackage{rotating}
\usepackage[normalem]{ulem}
\usepackage{capt-of}
\usepackage{hyperref}
\author{Prologos Project}
\date{2026-03-16}
\title{Tensor-Aware Propagator Networks with Unitary Evolution\\\medskip
\large A Research Program for Quantum-Like Computation in Prologos}
\hypersetup{
 pdfauthor={Prologos Project},
 pdftitle={Tensor-Aware Propagator Networks with Unitary Evolution},
 pdfkeywords={},
 pdfsubject={},
 pdfcreator={Emacs 30.2 (Org mode 9.7.39)}, 
 pdflang={English}}
\begin{document}

\maketitle
\tableofcontents

\section{1. Introduction and Motivation}
\label{sec:orgd45c75e}

Prologos's propagator network is a persistent, immutable computation engine where
cells hold lattice values and propagators are monotone functions that refine
information toward a fixed point. The Beyond Prolog research program (§6) proposed
extending this architecture with semiring-parameterized logic programming, where a
\texttt{DistCell} --- a cell holding a mapping from values to semiring weights --- enables
probabilistic, tropical, and fuzzy reasoning within the same infrastructure.

This document proposes the next step: extending beyond \texttt{DistCell} to \textbf{tensor-aware}
cells that hold joint amplitude distributions over multiple variables, enabling
quantum-like computation. The motivation is threefold.

\textbf{First, the modeling opportunity.} A substantial and growing body of research
applies the mathematical formalism of quantum mechanics --- superposition,
interference, entanglement, projection --- to model social, economic, and cognitive
phenomena. Quantum decision theory (Yukalov \& Sornette), quantum game theory
(Eisert, Wilkens, Lewenstein), quantum cognition (Busemeyer \& Bruza), and quantum
finance (Haven \& Khrennikov) all demonstrate that Hilbert space probability is a
more natural framework than classical (Kolmogorovian) probability for systems
exhibiting contextuality, order effects, and interference. A companion document
surveys this field in detail. These models require the ability to reason over
complex-valued amplitude distributions and tensor product state spaces --- exactly
what a tensor-aware propagator network would provide.

\textbf{Second, the architectural fit.} Prologos's existing infrastructure ---
persistent propagator networks, the Lattice trait hierarchy, ATMS for
hypothetical reasoning, stratified evaluation, and Galois connections for
cross-domain bridging --- provides a uniquely suitable substrate for this
extension. The critical insight, developed below, is that unitary evolution
(the quantum analog of "computation between observations") can be recovered
through the same \textbf{Layered Recovery Principle} that already governs the
interaction between monotonic propagation, non-monotonic ATMS control, and
stratified negation.

\textbf{Third, the categorical coherence.} The Beyond Prolog document (§7-8) established
that Prologos's logic engine should be parameterized by a quantale rather than
merely a lattice. The same quantale-enriched categorical framework that unifies
Boolean, probabilistic, tropical, and fuzzy logic programming also subsumes
quantum-like computation --- one needs only to instantiate the quantale with
complex-valued weights and non-commutative tensor products. The extension proposed
here is not an ad hoc addition but a natural instantiation of an already-designed
generalization.
\subsection{1.1 From DistCell to TensorCell}
\label{sec:org93a2973}

The \texttt{DistCell} design (Beyond Prolog §6.7, stage 4) holds a \texttt{Map Value Weight}
where merge applies pointwise semiring addition. This is sufficient for
\emph{classical} probabilistic reasoning --- each value carries an independent weight,
and the distribution is always separable across variables.

Quantum-like reasoning requires more. When two variables are \emph{entangled}, the
joint distribution over both variables cannot be factored into independent
marginals. The state of the composite system is described by a vector in the
\emph{tensor product} of the individual Hilbert spaces, and the weights are
complex-valued amplitudes (not real-valued probabilities). Interference arises
because amplitudes add before being squared to produce probabilities --- the
cross-terms in |a + b|² = |a|² + |b|² + 2Re(a*b) are the interference effects
that classical probability cannot represent.

A \texttt{TensorCell} holds a \texttt{Map (Tuple Value ...) Complex} --- a joint amplitude
distribution over multiple variables simultaneously. Its merge function applies
pointwise addition of complex amplitudes (superposition). Measurement (observation)
projects the state onto a subspace and applies the Born rule to convert amplitudes
to probabilities. The result is a cell type that can represent superposition,
interference, and entanglement within the same propagator infrastructure that
already handles type inference, session type checking, and relational logic.
\subsection{1.2 Document Structure}
\label{sec:orgf2be465}

This document proceeds as follows. Section 2 establishes the Layered Recovery
Principle as the architectural foundation, showing how it is already demonstrated
in three subsystems and how it extends to unitary evolution. Section 3 develops
the mathematical foundations --- from semirings through quantales to dagger
compact closed categories. Section 4 presents the TensorCell design in detail.
Section 5 shows how unitary evolution is recovered as a layered stratum. Section 6
develops the categorical architecture. Section 7 provides the implementation
roadmap. Section 8 connects to quantum social and economic models. Section 9
identifies research questions and open problems. Section 10 collects references.
\section{2. Background: The Layered Recovery Principle}
\label{sec:orgd4d18cc}

The central architectural insight of Prologos's logic engine is that
non-monotonic behavior need not be \emph{primitive}. Instead, it can be
\textbf{recovered} through carefully layered strata, each monotonic within itself,
with well-defined transitions between layers. We call this the \textbf{Layered
Recovery Principle}.
\subsection{2.1 The Three-Layer Architecture}
\label{sec:orgb22d394}

The logic engine design document established a three-layer architecture:

\begin{verbatim}
Layer 3: STRATIFICATION (controlled non-monotonicity)
    ↓
Layer 2: ATMS / TRUTH MAINTENANCE (non-monotonic control)
    ↓
Layer 1: PROPAGATOR NETWORK (monotonic data plane)
\end{verbatim}

\textbf{Layer 1} is the monotonic foundation. Cells hold lattice values; propagators
are monotone functions; the network runs to quiescence (fixed point) via
iterated propagation. Information only grows. The \texttt{run-to-quiescence} function
in \texttt{propagator.rkt} implements this as a worklist algorithm over the persistent
CHAMP-backed network. The BSP (Bulk Synchronous Parallel) variant
\texttt{run-to-quiescence-bsp} enables parallel execution via Jacobi iteration.

\textbf{Layer 2} introduces non-monotonic control \emph{without breaking monotonicity of
the data plane}. The ATMS (\texttt{atms.rkt}) wraps the propagator network with
assumptions, supported values, and nogoods. Each value carries a \emph{support set}
--- the assumptions that justify it. Multiple contingent values coexist in TMS
cells, and worldview switching selects which assumptions are currently believed.
The \texttt{atms-amb} operation creates mutually exclusive alternatives, and
dependency-directed backtracking identifies minimal inconsistent assumption sets.
Crucially, the propagator network underneath remains monotonic; the ATMS provides
the \emph{appearance} of non-monotonicity through contingent truth.

\textbf{Layer 3} adds stratified evaluation for negation-as-failure. The \texttt{stratify.rkt}
module computes SCCs (Strongly Connected Components) via Tarjan's algorithm and
ensures that no negative dependency cycle exists. The \texttt{stratified-eval.rkt}
orchestrator evaluates strata bottom-up, ensuring that each stratum reaches its
fixed point before the next stratum consults it via negation. Programs without
negation pay zero overhead.
\subsection{2.2 Demonstration in Three Subsystems}
\label{sec:orgbe65fef}

The Layered Recovery Principle is not merely theoretical. It is demonstrated in
three implemented subsystems:

\textbf{Type inference.} The elaborator network (\texttt{elaborator-network.rkt}) performs type
inference via monotonic propagation over a type lattice (\texttt{type-lattice.rkt}).
Metavariables become cells; unification constraints become bidirectional
propagators. When multiple type-checking strategies are possible (e.g., Church
fold vs. Scott fold for algebraic types), the speculative elaboration system
(\texttt{elab-speculation.rkt}) uses ATMS-backed fork-try-commit: each alternative forks
the network (O(1) persistent copy), applies elaboration, and checks for
contradiction. The first non-contradictory branch wins. Cross-domain propagation
via Galois connections bridges the type lattice to the multiplicity lattice (for
QTT) --- \texttt{net-add-cross-domain-propagator} in \texttt{propagator.rkt} creates
bidirectional α/γ propagators.

\textbf{Session types.} The session type checker (\texttt{session-propagators.rkt}) uses the
propagator network with a session lattice (\texttt{session-lattice.rkt}). Each channel
endpoint has a session cell holding the remaining protocol. Send/receive/choice
operations decompose protocols via propagators. Duality propagators enforce that
dual endpoints have compatible sessions. Protocol completeness checking detects
incomplete protocols (potential deadlocks). The session runtime
(\texttt{session-runtime.rkt}) executes session-typed processes over the persistent
network, with each channel advancement creating fresh cells rather than mutating.

\textbf{Relational logic.} The relation solver (\texttt{relations.rkt}) uses propagation for
constraint solving within a stratum, ATMS for choice points (\texttt{atms-amb} for
alternative clause selection), and stratified evaluation for negation-as-failure.
Tabling (\texttt{tabling.rkt}) provides SLG-style memoization with lattice answer modes,
and the persistent table store enables speculative evaluation without corruption.

In each case, the pattern is identical: a monotonic core (Layer 1) handles the
deterministic computation; a non-monotonic control layer (Layer 2) handles choice
and speculation; a stratified coordination layer (Layer 3) handles operations that
require prior layers to stabilize. The Galois connection machinery provides
well-typed transitions between different information domains.
\subsection{2.3 Effect Ordering as Interacting Quantales}
\label{sec:org915ab6b}

The effect ordering system provides a deeper perspective on layered recovery.
The system can be understood as two quantale structures interacting over quantale
morphisms:

\begin{itemize}
\item The \emph{information quantale} Q\textsubscript{I} tracks the lattice of partial information in
cells. Its join (\bigvee) merges information from alternative derivation paths;
its tensor (\(\otimes\)) combines information within a single derivation.
\item The \emph{control quantale} Q\textsubscript{C} tracks the ordering of effects --- which operations
must complete before others may begin. Its join represents nondeterministic
choice (ATMS alternatives); its tensor represents sequential composition
(stratum ordering).
\end{itemize}

A \emph{quantale morphism} φ: Q\textsubscript{C} → Q\textsubscript{I} maps control decisions to information updates.
When a stratum boundary is crossed (a control event), the morphism triggers the
appropriate information-level response (evaluating the next stratum's propagators
against the now-stable results of the previous stratum).

This two-quantale perspective reveals why the architecture is both principled and
extensible: adding a new kind of non-monotonic behavior (unitary evolution,
measurement) amounts to adding a new stratum to Q\textsubscript{C} and a new morphism to Q\textsubscript{I}.
The existing machinery handles the coordination.
\subsection{2.4 Why the Principle Extends to Unitary Evolution}
\label{sec:org9adf7a0}

Unitary evolution --- the quantum analog of "computation between observations"
--- is often presented as fundamentally incompatible with monotonic propagation.
The argument: propagators are monotonic (information only grows), but unitary
operators \emph{rotate} information without growing it. A rotation is
information-preserving but not information-increasing; it does not respect the
lattice ordering.

This argument is correct at the level of a single stratum. But the Layered
Recovery Principle dissolves it. The resolution:

\begin{enumerate}
\item \textbf{Unitary evolution is a stratum, not a propagator.} Just as negation-as-failure
is not a monotonic operation within a stratum but is \emph{recovered} by evaluating
strata in the right order, unitary evolution is not a monotonic operation within
the information lattice but can be \emph{recovered} by introducing a unitary stratum
between measurement points.

\item \textbf{Measurement is the stratum boundary.} A measurement (observation, projection)
is the event that triggers the transition from the unitary stratum to the
information stratum. Before measurement, the system evolves unitarily (rotating
the state vector). At measurement, the state is projected onto an eigenspace,
producing a classical probability distribution (via the Born rule) that \emph{is}
monotonic information.

\item \textbf{The Galois connection bridges the gap.} The transition from amplitude space
(complex-valued, pre-measurement) to probability space (real-valued,
post-measurement) is a Galois connection. The abstraction function α applies
the Born rule (|ψ|²); the concretization function γ embeds classical
probabilities as diagonal density matrices. This is structurally identical to
how \texttt{net-add-cross-domain-propagator} bridges the type lattice to the
multiplicity lattice.

\item \textbf{ATMS handles measurement outcomes.} Each possible measurement outcome is an
ATMS assumption. Different outcomes are mutually exclusive alternatives,
created by \texttt{atms-amb}. The nogood machinery handles contradictory outcomes.
This is structurally identical to how the elaborator handles speculative
type checking.
\end{enumerate}

The result: unitary evolution is \emph{recovered} within the existing architecture by
adding two new strata (unitary evolution and measurement) and one new Galois
connection (amplitudes ↔ probabilities), using the same mechanisms that already
handle type inference, session checking, and relational logic.
\section{3. Mathematical Foundations}
\label{sec:orga4b0a64}

\subsection{3.1 From Semirings to Quantales}
\label{sec:org4cdcd97}

The Beyond Prolog document (§6-7) established the progression from lattices to
semirings to quantales. We recap the key structures and extend them to the
complex-valued setting needed for quantum-like computation.

A \textbf{semiring} (S, ⊕, ⊗, \bar{0}, \bar{1}) has two binary operations:
addition ⊕ (with identity \bar{0}) and multiplication ⊗ (with identity \bar{1}).
Addition is commutative; multiplication distributes over addition. The Beyond
Prolog document showed how different semiring instantiations yield different modes
of logic programming:

\begin{center}
\begin{tabular}{llll}
Semiring & ⊕ & ⊗ & Question answered\\
\hline
Boolean & ∨ & ∧ & Is there any derivation?\\
Probability & + & × & What is the total probability?\\
Tropical & min & + & What is the minimum cost?\\
Counting & + & × & How many derivations exist?\\
Viterbi & max & × & What is the most probable path?\\
Provenance & ∪ & ⋈ & Which base facts were used?\\
\end{tabular}
\end{center}

A \textbf{quantale} (Q, ≤, ⊗) is a complete lattice equipped with an associative tensor
⊗ distributing over arbitrary joins. Quantales subsume both lattices (⊗ = ∧) and
semirings (≤ is the additive order). The critical distinction: join (\bigvee) merges
information from \emph{different} derivation paths; tensor (⊗) combines information
from the \emph{same} derivation.
\subsection{3.2 Complex Semirings and Amplitude Spaces}
\label{sec:orgc4bee4e}

For quantum-like computation, we instantiate the semiring with complex numbers:

\begin{center}
\begin{tabular}{llll}
Semiring & ⊕ & ⊗ & Question answered\\
\hline
Complex & + & × & What is the total amplitude?\\
\end{tabular}
\end{center}

The complex semiring (ℂ, +, ×, 0, 1) is the natural setting for quantum-like
reasoning. Key properties that distinguish it from the probability semiring:

\begin{itemize}
\item \textbf{Interference.} Complex addition can be constructive or destructive. If two
derivation paths contribute amplitudes a and b to the same conclusion, the
total amplitude is a + b, and the probability is |a + b|² = |a|² + |b|² +
2Re(a*b). The cross-term 2Re(a*b) is the interference term --- absent in
classical probability, where probabilities simply add.

\item \textbf{Phase.} Complex numbers carry phase (direction in the complex plane). Two
amplitudes with opposite phases cancel; two with aligned phases reinforce.
Phase is the mechanism by which quantum-like models produce non-classical
probability patterns.

\item \textbf{Unitarity.} Transformations that preserve the total probability (Σ|aᵢ|² = 1)
are unitary matrices. These are the "legal moves" between observations ---
they rotate the state vector without changing its length.
\end{itemize}

The amplitude space forms a lattice under the information ordering of
sub-probability distributions (§3.3), but the \emph{phase} information is invisible
to this ordering. This is why unitary evolution does not respect the lattice
order --- it changes phases (and hence interference patterns) without changing
the total information content.
\subsection{3.3 Tensor Products of Information Lattices}
\label{sec:orge930f81}

If cell A holds values from lattice L\textsubscript{A} and cell B from lattice L\textsubscript{B}, a constraint
relating both variables lives in the tensor product L\textsubscript{A} ⊗ L\textsubscript{B}. The tensor product
of two join-semilattices is the universal recipient of bihomomorphisms ---
functions that are lattice homomorphisms in each argument separately.

Arntzenius proved that the category \textbf{SemiLat} of join-semilattices forms a
closed symmetric monoidal category. The internal hom A ⊸ B is the semilattice
of homomorphisms ordered pointwise. This means linear lambda-calculus is a sound
typed language for defining semilattice homomorphisms --- a connection to
Prologos's linear types (QTT) that deserves further exploration.

For quantum-like computation, the relevant tensor product is that of Hilbert
spaces. If system A has basis \{|0⟩, |1⟩\} and system B has basis \{|0⟩, |1⟩\},
the composite system A⊗B has basis \{|00⟩, |01⟩, |10⟩, |11⟩\}. A state in
A⊗B is a vector in this 4-dimensional space:

\begin{verbatim}
|ψ⟩ = α|00⟩ + β|01⟩ + γ|10⟩ + δ|11⟩
\end{verbatim}

When the state \emph{cannot} be factored as |ψ\textsubscript{A}⟩ ⊗ |ψ\textsubscript{B}⟩, the systems are
\textbf{entangled} --- their joint state contains correlations that no product of
individual states can reproduce. The Bell state (|00⟩ + |11⟩)/√2 is the
canonical example: measuring A determines B's value with certainty, yet
neither A nor B individually has a definite value.

For the propagator network, this means a \texttt{TensorCell} covering variables A and B
holds a 4-entry map \{|00⟩↦α, |01⟩↦β, |10⟩↦γ, |11⟩↦δ\}, and the entanglement
manifests as the inability to decompose this into independent per-variable cells.

Tensor network contraction methods from physics (Biamonte et al., 2018) provide
efficient algorithms for evaluating multi-cell constraints through such tensor
product spaces. These methods represent constraint satisfaction as tensor network
contraction and have been shown to outperform state-of-the-art solvers on certain
problem classes (\#monotone 1-in-3SAT, \#Cubic-Vertex-Cover).
\subsection{3.4 Dagger Categories and Unitary Morphisms}
\label{sec:orgde40646}

Categorical quantum mechanics (Abramsky \& Coecke, 2004) axiomatizes quantum
theory in terms of dagger compact closed categories. The key structures:

A \textbf{dagger category} has a contravariant involution on morphisms: every morphism
f: A → B has an adjoint f†: B → A, with (f†)† = f and (g ∘ f)† = f† ∘ g†.
A \textbf{unitary} morphism satisfies u† ∘ u = id = u ∘ u†. Unitaries are the
"reversible" morphisms --- they transform without losing information.

A \textbf{dagger compact closed category} additionally has:
\begin{itemize}
\item A symmetric monoidal structure (⊗, I) for parallel composition.
\item Compact closure: every object A has a dual A*, with cups (η: I → A ⊗ A*) and
caps (ε: A* ⊗ A → I) satisfying the snake equations.
\item Compatibility: the dagger respects the monoidal structure.
\end{itemize}

The canonical example is \textbf{FdHilb}: finite-dimensional Hilbert spaces and linear
maps. Objects are Hilbert spaces; morphisms are linear maps; the dagger is the
adjoint (conjugate transpose); the tensor is the Hilbert space tensor product;
the dual A* is the conjugate Hilbert space.

For a propagator network, this structure maps as follows:
\begin{itemize}
\item \textbf{Objects} (Hilbert spaces) correspond to cell types --- the state spaces of
individual cells or tensor products of multiple cells.
\item \textbf{Morphisms} (linear maps) correspond to propagators --- transformations of
cell values.
\item \textbf{Unitary morphisms} correspond to unitary propagators --- reversible
transformations that preserve total probability.
\item \textbf{The dagger} gives the inverse/adjoint of every propagator, enabling
"reverse propagation."
\item \textbf{Compact closure} enables bidirectional constraint propagation (running
constraints forwards or backwards by "bending wires").
\end{itemize}

Selinger (2007) proved that the graphical language of dagger compact closed
categories (string diagrams) is complete for FdHilb: any equation that holds
for finite-dimensional Hilbert spaces can be proved by diagram manipulation.
This means a graphical representation of a tensor-aware propagator network
would have a \emph{complete} equational theory.
\subsection{3.5 Frobenius Algebras as Merge/Split Operations}
\label{sec:orgffd85d6}

A \textbf{Frobenius algebra} in a dagger compact closed category is an object equipped
with a multiplication (merge) μ: A ⊗ A → A and a comultiplication (split)
δ: A → A ⊗ A satisfying the Frobenius law: (μ ⊗ id) ∘ (id ⊗ δ) = δ ∘ μ =
(id ⊗ μ) ∘ (δ ⊗ id).

In FdHilb, dagger commutative Frobenius algebras correspond one-to-one with
orthonormal bases (Coecke, Pavlovic \& Vicary, 2012). The multiplication copies
basis elements: μ(|i⟩ ⊗ |j⟩) = δᵢⱼ|i⟩ (merge = agree-or-contradict). The
comultiplication broadcasts: δ(|i⟩) = |i⟩ ⊗ |i⟩ (split = copy). For a
propagator network, this provides the algebraic structure for merge points ---
where information from multiple sources combines.

The \textbf{spider theorem} states that any connected diagram built from a single
Frobenius algebra depends only on the number of inputs and outputs. A "spider"
with n inputs and m outputs has a canonical form regardless of internal wiring.
For propagator networks, this means the \emph{topology} of merging is irrelevant ---
only the \emph{arity} matters. Whether three propagators merge pairwise or all at
once produces the same result.

\textbf{Complementary observables} arise from two Frobenius algebras that interact via
a bialgebra structure. In quantum mechanics, position and momentum are
complementary: sharp knowledge of one implies maximal uncertainty about the
other. In a propagator network, this structure models two incompatible
constraint domains that interact in a controlled way --- reminiscent of
the cross-domain propagation via Galois connections that Prologos already
supports, but with a richer interaction structure.
\subsection{3.6 The CPM Construction: Pure to Mixed Propagation}
\label{sec:org9f017ad}

Not all quantum processes are pure (reversible, unitary). Interaction with an
environment introduces noise, decoherence, and information loss. The CPM
(Completely Positive Maps) construction (Selinger, 2007) provides the
categorical machinery:

Given a dagger compact closed category C (modeling pure processes), \textbf{CPM(C)} is
another dagger compact closed category modeling mixed processes. The functor
F: C → CPM(C) maps each morphism f to f* ⊗ f (where f* is the conjugate),
converting amplitudes to probabilities. Objects in CPM(C) are the same as in C,
but morphisms A → B in CPM(C) are morphisms A ⊗ A* → B ⊗ B* in C that are
"completely positive" --- they map positive operators to positive operators even
when tensored with an identity.

Coecke (2008) recast this in terms of \textbf{environment structures}: a designated
"discarding" effect ⊤\textsubscript{A}: A → I for each object, representing what is not
observed. Pure morphisms form a sub-category; mixed morphisms additionally
interact with the environment.

The CP* construction (Coecke, Heunen \& Kissinger) generalizes further: objects
are "abstract C*-algebras" (dagger Frobenius algebras), unifying quantum
channels and classical stochastic maps in a single category. Classical systems
are commutative C*-algebras; quantum systems are noncommutative ones.

For the propagator network, this provides the bridge from ideal propagation
(fully tracked, reversible) to realistic propagation (partial observation,
information loss):
\begin{itemize}
\item \textbf{Pure propagators} are unitary morphisms in C --- fully reversible, no
information loss. These govern evolution between measurements.
\item \textbf{Mixed propagators} are morphisms in CPM(C) --- they may lose information,
modeling noisy channels, partial observations, or environmental decoherence.
\item \textbf{Discarding maps} are what happens when you trace out part of a tensor
product --- i.e., marginalization. This connects directly to the
\texttt{marginalize} operation proposed in Beyond Prolog §6.7.
\item The CP* unification means that \textbf{classical cells} (DistCells with real weights)
and \textbf{quantum cells} (TensorCells with complex amplitudes) coexist in the same
network, with classical cells being the commutative special case.
\end{itemize}
\section{4. The TensorCell Design}
\label{sec:orgec0e1c6}

\subsection{4.1 Cell Hierarchy}
\label{sec:org594cd8d}

The extension from classical to quantum-like cells proceeds through four levels,
each subsuming the previous:
\subsubsection{Level 1: DistCell --- Classical Distributions}
\label{sec:orga468ce9}

\begin{verbatim}
;; A cell holding a classical probability distribution
;; Map from values to real-valued semiring weights
type DistCell {V} := Map V Real
\end{verbatim}

\texttt{DistCell} holds a mapping from values to real-valued weights. Merge is
pointwise semiring addition. This is the design from Beyond Prolog §6.7,
sufficient for probabilistic databases, Bayesian reasoning, tropical
optimization, and counting.

\textbf{Information ordering:} d₁ ≤ d₂ iff d₁(v) ≤ d₂(v) for all v (pointwise).
Bottom is the zero distribution. This forms a complete lattice suitable for
propagation.
\subsubsection{Level 2: AmplitudeCell --- Single-System Quantum States}
\label{sec:org3b0da11}

\begin{verbatim}
;; A cell holding a quantum-like amplitude distribution
;; Map from values to complex-valued weights
type AmplitudeCell {V} := Map V Complex
\end{verbatim}

\texttt{AmplitudeCell} extends \texttt{DistCell} with complex-valued weights. Merge is
pointwise complex addition (superposition). The probability of a value v is
\begin{center}
\begin{tabular}{ll}
amplitude(v) & ² (Born rule).\\
\end{tabular}
\end{center}

\textbf{Key new capability:} Interference. When two derivation paths contribute
amplitudes a and b to the same value, the total probability is |a+b|², not
\begin{center}
\begin{tabular}{llll}
a & ²+ & b & ². Destructive interference (a ≈ -b) can make a classically probable\\
\end{tabular}
\end{center}
outcome improbable; constructive interference can amplify rare outcomes.

This level captures most of the quantum cognition literature: order effects,
conjunction fallacy, disjunction effect, preference reversal. These are all
single-system phenomena where one agent's mental state is in superposition.
\subsubsection{Level 3: TensorCell --- Multi-System Entangled States}
\label{sec:orgb62444c}

\begin{verbatim}
;; A cell holding a joint amplitude distribution over multiple variables
;; Map from tuples of values to complex weights
;; This is a state vector in a tensor-product Hilbert space
type TensorCell {Vs : Vec n Type} := Map (Tuple Vs) Complex
\end{verbatim}

\texttt{TensorCell} holds a joint amplitude distribution over \emph{multiple} variables
simultaneously. The key space is the tensor product of individual variable
spaces. A state in the tensor product may be \textbf{entangled} --- non-factorable
into independent per-variable states.

\textbf{Key new capability:} Entanglement. Two agents' belief states, two financial
instruments' price states, two voters' preference states --- any pair of
variables can have correlations that no product of marginals can reproduce.

\textbf{Dependent typing:} The \texttt{Vs : Vec n Type} parameter tracks the dimensionality
at the type level. A \texttt{TensorCell \{[Bool Bool]\}} is a 4-dimensional state
space; a \texttt{TensorCell \{[Bool Bool Bool]\}} is 8-dimensional. The dependent type
ensures dimensional consistency at compile time.
\subsubsection{Level 4: DensityCell --- Mixed States}
\label{sec:org51428c5}

\begin{verbatim}
;; A cell holding a density matrix (mixed state)
;; Map from pairs of basis elements to complex weights
;; ρ = Σᵢ pᵢ |ψᵢ⟩⟨ψᵢ|
type DensityCell {V} := Map (Pair V V) Complex
\end{verbatim}

\texttt{DensityCell} holds a density matrix --- the most general representation of a
quantum(-like) state, including both pure states and statistical mixtures.
Diagonal entries ρ(v,v) are probabilities; off-diagonal entries ρ(v,w) are
coherences (encoding interference capability).

\textbf{Key new capability:} Mixed states. When you don't know \emph{which} pure state a
system is in, but only a probability distribution over pure states, the density
matrix encodes this uncertainty. Decoherence (loss of quantum effects through
environmental interaction) manifests as decay of off-diagonal entries toward
zero, leaving a classical probability distribution on the diagonal.

This level enables the open quantum systems approach used in cognitive modeling
(Asano, Khrennikov et al.) where decision-making is modeled as decoherence of
a mental state --- the GKSL master equation governs the evolution.
\subsection{4.2 Merge Semantics}
\label{sec:orgb2e54eb}

Each cell level requires a merge function compatible with the Lattice trait:
\subsubsection{DistCell merge}
\label{sec:org868096b}
Pointwise real addition. d₁ ⊕ d₂ maps each value v to d₁(v) + d₂(v).
Monotonic under the pointwise ordering.
\subsubsection{AmplitudeCell merge}
\label{sec:orga7cc515}
Pointwise complex addition (superposition). a₁ ⊕ a₂ maps each value v to
a₁(v) + a₂(v). This produces interference: the merged amplitude may be larger
(constructive) or smaller (destructive) than either input.

The information ordering for amplitudes uses the \emph{sub-distribution} lattice:
a₁ ≤ a₂ iff the support of a₁ is a subset of the support of a₂ and the total
amplitude magnitude of a₁ is at most that of a₂. Under this ordering, merge
is monotonic: adding more amplitude information (more derivation paths) only
increases the total.
\subsubsection{TensorCell merge}
\label{sec:org94fa38d}
Pointwise complex addition over the tensor product basis. The merge of two
TensorCells combines amplitude contributions to each basis element. If two
derivation paths contribute amplitudes to the same joint outcome |v₁,v₂,\ldots{}⟩,
they superpose.

The ordering is the pointwise sub-distribution ordering over the tensor product
basis. Since the tensor product basis is finite (for finite-dimensional systems),
this is a complete lattice.
\subsubsection{DensityCell merge}
\label{sec:org530aa3b}
Convex combination or pointwise addition of density matrices, depending on the
semantics. For combining evidence from independent sources, pointwise addition
(followed by renormalization) is appropriate. For representing classical mixture
of quantum states, convex combination (ρ = p₁ρ₁ + p₂ρ₂) is appropriate.
\subsection{4.3 Tensor Product Structure}
\label{sec:org4358efb}

The tensor product of two cells A and B produces a new cell whose state space
is the Cartesian product of basis elements with amplitude given by the product:

\begin{verbatim}
(A ⊗ B)(|i,j⟩) = A(|i⟩) × B(|j⟩)
\end{verbatim}

This is the "no entanglement" case --- a product state. Entanglement arises when
the joint state cannot be expressed as a product.

Operations on tensor products:

\textbf{Partial trace (marginalization).} Summing over one subsystem's basis elements:
tr\textsubscript{B}(ρ\textsubscript{AB}) maps the joint state to a reduced state on A alone. This is the
quantum analog of marginalization in classical probability. For a TensorCell
over variables (A,B), tracing out B produces an AmplitudeCell over A.

\textbf{Tensor contraction.} More general than partial trace --- contracts a pair of
indices, potentially connecting two tensor cells. This is the fundamental
operation of tensor network methods and maps directly to propagator wiring.

\textbf{Schmidt decomposition.} Any bipartite pure state |ψ\textsubscript{AB}⟩ can be written as
Σᵢ λᵢ|aᵢ⟩|bᵢ⟩ where λᵢ ≥ 0 and \{|aᵢ⟩\}, \{|bᵢ⟩\} are orthonormal. The
Schmidt rank measures entanglement: rank 1 = separable; rank > 1 = entangled.
For a TensorCell, the Schmidt decomposition provides a compressed representation
and a measure of entanglement degree.
\subsection{4.4 Trait Instances}
\label{sec:org0e2de2e}

The cell hierarchy integrates with Prologos's trait system:

\begin{verbatim}
;; All cell types implement Lattice via their merge semantics
impl Lattice (AmplitudeCell V) where (Eq V)
  bot  := empty-map
  join := pointwise-complex-add
  leq  := pointwise-magnitude-leq

;; TensorCell is a Lattice (pointwise over tensor product basis)
impl Lattice (TensorCell Vs) where (Eq (Tuple Vs))
  bot  := empty-map
  join := pointwise-complex-add
  leq  := pointwise-magnitude-leq

;; Quantale instance: tensor = pointwise complex multiplication
impl Quantale (AmplitudeCell V) where (Eq V)
  one := uniform-distribution
  mul := pointwise-complex-mul

;; Widenable for approximation of high-dimensional spaces
impl Widenable (TensorCell Vs)
  widen  := truncate-low-amplitude-entries
  narrow := restore-from-exact
\end{verbatim}

The \texttt{Widenable} instance is important for practical computation: high-dimensional
tensor product spaces grow exponentially with the number of variables. Widening
by truncating entries with amplitude below a threshold provides a controlled
approximation, analogous to how interval widening jumps to ±∞ to ensure
convergence.
\section{5. Unitary Evolution as a Layered Stratum}
\label{sec:orgfe1e6a8}

\subsection{5.1 The Monotonicity Tension}
\label{sec:org0fe4926}

Standard propagator semantics require monotonicity: if a cell's value is v and
a propagator writes v', the cell's new value is join(v, v'), which is ≥ v in
the lattice ordering. Information never decreases.

Unitary evolution does not satisfy this. A unitary operator U applied to a state
\begin{center}
\begin{tabular}{ll}
ψ⟩ produces U & ψ⟩, which has the same norm (total probability) but potentially\\
\end{tabular}
\end{center}
different amplitudes for each basis element. Some amplitudes increase; others
decrease. The state is \emph{rotated}, not \emph{refined}.

This is directly analogous to the monotonicity tension with negation-as-failure.
Negation inverts: if a fact is derived, its negation should retract. This violates
monotonicity within a single stratum. The solution (stratification) is to evaluate
the positive part to completion \emph{first}, then apply negation against the stable
result.
\subsection{5.2 The Resolution: A Four-Layer Architecture}
\label{sec:org00e10e8}

The tensor-aware propagator network extends the three-layer architecture with
a fourth layer:

\begin{verbatim}
Layer 4: MEASUREMENT (projection, Born rule → classical probabilities)
    ↓
Layer 3: UNITARY EVOLUTION (rotations in amplitude space)
    ↓
Layer 2: ATMS / TRUTH MAINTENANCE (non-monotonic control)
    ↓
Layer 1: MONOTONIC PROPAGATION (lattice joins, constraint refinement)
\end{verbatim}

The execution model:

\begin{enumerate}
\item \textbf{Layer 1 (Monotonic propagation).} Classical constraint propagation runs to
quiescence. Cells hold lattice values; propagators refine them. This handles
all the "classical" parts of the computation: type checking, mode inference,
domain constraints.

\item \textbf{Layer 2 (ATMS control).} Choice points are introduced via \texttt{atms-amb}. Each
alternative represents a branch of nondeterministic search. Nogoods prune
impossible branches. This handles relational search, speculative elaboration,
and (for quantum models) the selection of which unitary evolution to apply.

\item \textbf{Layer 3 (Unitary evolution).} After the classical constraints have stabilized
(Layer 1 quiescent) and the nondeterministic choices have been made (Layer 2),
unitary operators are applied to TensorCells. This is a \emph{single pass} of
unitary evolution --- no iteration, no fixed point, just matrix application.
Multiple unitary operators compose (U₂ ∘ U₁) to form the net evolution.

\item \textbf{Layer 4 (Measurement).} After unitary evolution completes, measurement
projectors are applied. Each measurement outcome is an ATMS assumption.
The Born rule converts amplitudes to probabilities. The resulting classical
probabilities feed back into Layer 1 for further constraint propagation.
\end{enumerate}

The cycle may repeat: Layer 1 propagates the measurement results, Layer 2 makes
new choices, Layer 3 applies new unitaries, Layer 4 measures again. This
corresponds to the standard quantum computing paradigm of interleaved evolution
and measurement.
\subsection{5.3 Unitary Propagators}
\label{sec:orgc21177a}

A \textbf{unitary propagator} is a new propagator type that applies a unitary matrix to
a TensorCell. Unlike standard propagators (which are monotone functions), unitary
propagators are \emph{isometries} --- they preserve the norm but rotate the state.

\begin{verbatim}
;; A unitary propagator applies a unitary matrix U to a TensorCell
;; U : Matrix Complex n n where U†U = I
spec apply-unitary : {n : Nat} -> Matrix Complex n n -> TensorCell n -> TensorCell n
defn apply-unitary [U cell] (matrix-mul U cell)

;; Composition of unitaries is matrix multiplication
spec compose-unitary : Matrix Complex n n -> Matrix Complex n n -> Matrix Complex n n
defn compose-unitary [U1 U2] (matrix-mul U1 U2)

;; The dagger (adjoint) is conjugate transpose
spec dagger : Matrix Complex n n -> Matrix Complex n n
defn dagger [U] (matrix-conjugate-transpose U)

;; Tensor product of unitaries for parallel evolution
spec tensor-unitary : Matrix Complex m m -> Matrix Complex n n
                   -> Matrix Complex [m * n] [m * n]
defn tensor-unitary [U1 U2] (matrix-kronecker U1 U2)
\end{verbatim}

Standard unitary gates from quantum computing map directly:

\begin{verbatim}
;; Hadamard: creates superposition
def H : Matrix Complex 2 2 := matrix-from-rows
  '[[0.7071 0.7071] [0.7071 -0.7071]]

;; Phase gate: adds relative phase
def S : Matrix Complex 2 2 := matrix-from-rows
  '[[1.0 0.0] [0.0 (complex 0.0 1.0)]]

;; CNOT: entangling gate
def CNOT : Matrix Complex 4 4 := matrix-from-rows
  '[[1 0 0 0] [0 1 0 0] [0 0 0 1] [0 0 1 0]]
\end{verbatim}

For quantum social science models, the unitary operators would encode:
\begin{itemize}
\item \textbf{Preference rotation:} How a decision-maker's belief state evolves over time.
\item \textbf{Social influence:} How interaction with another agent rotates one's state.
\item \textbf{Entangling operations:} How two agents become correlated through interaction.
\end{itemize}
\subsection{5.4 Measurement as Projection}
\label{sec:orga39f79d}

Measurement in quantum mechanics is projection onto an eigenspace, followed by
normalization (the Born rule). In the propagator network:

\begin{verbatim}
;; Measurement projects onto a subspace and computes probability
spec measure : {V : Type} -> AmplitudeCell V -> V -> Pair Real (AmplitudeCell V)
defn measure [cell value]
  let amplitude := (map-lookup cell value)
  let probability := (complex-magnitude-squared amplitude)
  let projected := (singleton-map value (complex 1.0 0.0))
  (pair probability projected)

;; Marginalization: trace out one subsystem
spec marginalize : {A B : Type} -> TensorCell [A B] -> AmplitudeCell A
defn marginalize [joint]
  ;; Sum |amplitude(a,b)|² over all b for each a
  (fold-map joint empty-map
    (fn [(pair (tuple a b) amp) acc]
      (map-update acc a (complex 0.0 0.0)
        (fn [current] (complex-add current
          (complex-mul (complex-conjugate amp) amp))))))
\end{verbatim}

The measurement operation is the Galois connection between the unitary and
classical layers:

\begin{itemize}
\item \textbf{Abstraction α (Born rule):} Maps TensorCell amplitudes to DistCell
probabilities via |ψ(v)|². This is a many-to-one map: many different amplitude
distributions can produce the same probability distribution.
\item \textbf{Concretization γ (embedding):} Maps DistCell probabilities to DensityCell
density matrices via the diagonal embedding ρ(v,v) = p(v), ρ(v,w) = 0 for
v ≠ w. This embeds classical distributions as "fully decohered" quantum states.
\item \textbf{Adjunction:} α(ψ) ≤ d iff ψ ≤ γ(d) holds under the appropriate information
orderings.
\end{itemize}
\subsection{5.5 Integration with ATMS}
\label{sec:org579be9f}

When a measurement is performed, each possible outcome becomes an ATMS assumption:

\begin{verbatim}
;; Pseudo-code for measurement integration
(define (measure-tensor-cell atms cell-id basis-values)
  ;; Create mutually exclusive assumptions for each outcome
  (define-values (atms* hyps)
    (atms-amb atms (map symbol->string basis-values)))

  ;; For each outcome, record the post-measurement state
  (for/fold ([a atms*]) ([hyp hyps] [val basis-values])
    (define amplitude (tensor-cell-lookup cell-id val))
    (define probability (complex-magnitude-squared amplitude))
    ;; Write the collapsed state to the cell under this assumption
    (atms-write-cell a cell-id
      (supported-value (singleton-tensor val) (hasheq hyp #t)))))
\end{verbatim}

This is structurally identical to how \texttt{speculation-begin} in \texttt{elab-speculation.rkt}
creates mutually exclusive branches for speculative type checking. The ATMS
machinery handles:
\begin{itemize}
\item \textbf{Mutual exclusion:} Only one measurement outcome can be actual.
\item \textbf{Dependency tracking:} Downstream conclusions carry the measurement assumption
in their support sets.
\item \textbf{Nogood detection:} If a measurement outcome leads to contradiction, it is
recorded as a nogood and pruned.
\item \textbf{Worldview exploration:} \texttt{atms-solve-all} explores all consistent worldviews,
each corresponding to a different measurement outcome.
\end{itemize}
\section{6. Categorical Architecture}
\label{sec:org44111c4}

\subsection{6.1 The Enrichment Base}
\label{sec:org09b68c5}

Prologos's propagator network is naturally a category enriched over a lattice: the
hom-object between two cells is not a set of morphisms but a lattice element
measuring "how much" information flows between them. The Beyond Prolog document
(§7.1) identified this as enrichment over a quantale --- a complete lattice with
an associative tensor distributing over arbitrary joins.

For tensor-aware propagation, the enrichment base must be extended. The relevant
quantale is \textbf{non-commutative}: the tensor product of two amplitude spaces is
sensitive to ordering (measuring A then B differs from measuring B then A). This
connects to Mulvey's (1986) original motivation for quantales as "quantum locales"
--- non-commutative generalizations of the lattice of open sets of a topological
space.

The spectrum Max(A) of a C*-algebra A is a quantale whose points correspond to
irreducible representations (Kruml, Pelletier, Resende, 2003). For
finite-dimensional quantum systems, this is the quantale of projection lattices.
The connection to Prologos: each cell type defines a C*-algebra (the algebra of
observables on that cell's state space), and the quantale of projections provides
the information ordering.

Stubbe (2013) generalized from single quantales to \textbf{quantaloids} --- many-object
quantales. A quantaloid-enriched category has hom-objects drawn from potentially
different quantales for different pairs of objects. This is exactly what a
heterogeneous propagator network requires: type cells are enriched over the type
lattice, session cells over the session lattice, amplitude cells over the complex
amplitude quantale, and cross-domain propagators correspond to quantaloid
morphisms between these enrichment bases.
\subsection{6.2 Dagger Compact Closed Structure}
\label{sec:org0b6d65c}

The tensor-aware propagator network lives in a dagger compact closed category
where:

\begin{itemize}
\item \textbf{Objects} are cell state spaces: type lattice, session lattice, amplitude
space, tensor product space, density matrix space.
\item \textbf{Morphisms} are propagators: monotone functions (Layer 1), ATMS operations
(Layer 2), unitary operators (Layer 3), projection/measurement (Layer 4).
\item \textbf{The dagger} maps each propagator to its adjoint. For unitary propagators,
this is the inverse. For measurement propagators, the dagger gives the
"preparation" (embedding a classical outcome back into amplitude space).
\item \textbf{The tensor product} combines cell state spaces for multi-variable constraints.
Two cells A and B can be combined into a TensorCell A⊗B.
\item \textbf{Compact closure} provides cups and caps --- the ability to "bend wires" in
string diagrams. This corresponds to running a constraint bidirectionally,
which is already a fundamental capability of propagator networks (a propagator
that refines cell A based on cell B also implicitly constrains B based on A).
\end{itemize}

The string diagram calculus provides a complete graphical language for reasoning
about the network. Every valid diagram manipulation corresponds to a valid
network transformation, and conversely.
\subsection{6.3 Frobenius Structure at Merge Points}
\label{sec:orgab52be7}

When multiple propagators write to the same cell, their contributions are merged
via the lattice join. For TensorCells, this merge has Frobenius algebra structure:

\begin{itemize}
\item \textbf{Multiplication} (merge): When two amplitude contributions arrive at the same
cell, they are added (superposed). This is the Frobenius multiplication
μ: A ⊗ A → A.
\item \textbf{Comultiplication} (broadcast): When a cell's value is read by multiple
propagators, it is copied to each. This is the Frobenius comultiplication
δ: A → A ⊗ A.
\end{itemize}

The \textbf{spider theorem} guarantees that any connected network of merges and broadcasts
built from a single Frobenius algebra reduces to a single "spider" determined
only by its number of inputs and outputs. For the propagator network, this means:
\emph{the order in which propagators write to a cell does not matter} --- only the
total set of contributions matters. This is a generalization of the confluence
property that makes propagator networks deterministic despite nondeterministic
scheduling.

When two \textbf{complementary} Frobenius algebras interact (corresponding to two
incompatible constraint domains), the bialgebra structure controls their
interaction. This generalizes the cross-domain propagation via Galois connections
to a richer algebraic interaction structure.
\subsection{6.4 ZX-Calculus as Rewrite System}
\label{sec:org888961e}

The ZX-calculus (Coecke \& Duncan, 2008; van de Wetering, 2020) is a complete
graphical rewrite system for qubit quantum mechanics, built from two families of
spiders (Z-type and X-type) interacting via complementary Frobenius algebras.

The rewrite rules map to propagator network operations:

\begin{center}
\begin{tabular}{ll}
ZX Rule & Propagator Network Operation\\
\hline
Spider fusion & Merging two propagators into one\\
Bialgebra rule & Cross-domain interaction\\
Copy rule & Broadcasting through complementary domain\\
Color change & Basis transformation (Hadamard)\\
Hopf rule & Strong complementarity simplification\\
\end{tabular}
\end{center}

\textbf{Completeness} means any valid equivalence between network configurations can be
discovered by local rewriting. This opens the possibility of \textbf{automated
optimization}: given a tensor-aware propagator network, apply ZX-calculus rewrites
to simplify it before execution. The circuit extraction problem (converting
optimized ZX-diagrams back to executable form) parallels the problem of extracting
an efficient execution plan from an optimized propagator network.

The PyZX library provides an open-source implementation of ZX-calculus rewriting
that could serve as a reference for a Prologos-native optimizer.
\subsection{6.5 The CPM Layer}
\label{sec:orgcc191ca}

The CPM construction provides the categorical machinery for the measurement
layer (Layer 4 of the architecture). The key insight: measurement is not a
separate operation bolted onto the propagator network; it is a \emph{functor}
from the pure category (TensorCells with unitary propagators) to the mixed
category (DensityCells with completely positive maps).

The discarding maps ⊤\textsubscript{A}: A → I in the environment structure correspond to
partial trace (marginalization). When a multi-variable TensorCell is observed
on a subset of its variables, the unobserved variables are "discarded" (traced
out), producing a reduced state on the observed variables.

This provides the formal foundation for the Galois connection between the
unitary layer (Layer 3) and the measurement layer (Layer 4):

\begin{verbatim}
Pure (Layer 3)        CPM functor        Mixed (Layer 4)
─────────────────     ──────────────>    ──────────────────
TensorCell            ─────>             DensityCell
Unitary propagator    ─────>             CP map
Tensor product        ─────>             Tensor product
Adjoint (†)           ─────>             Adjoint (†)
\end{verbatim}

Classical cells (DistCells) embed into the mixed category as commutative
C*-algebras (diagonal density matrices). This ensures backward compatibility:
all existing propagator network behavior is a special case of the tensor-aware
extension.
\section{7. Implementation Roadmap}
\label{sec:org08f531e}

The implementation proceeds in eight phases, each delivering independently
testable functionality. The phases build on the existing infrastructure
(propagator network, lattice traits, ATMS, Galois connections) and extend the
Beyond Prolog probabilistic LP roadmap.
\subsection{7.1 Phase T1: Complex Number Type and Semiring}
\label{sec:org3ac31b2}

\textbf{Deliverables:}
\begin{itemize}
\item Complex number type with arithmetic operations
\item \texttt{Semiring} trait instance for Complex
\item \texttt{Lattice} instance for complex-valued distributions (pointwise ordering)
\end{itemize}

\textbf{Dependencies:} None beyond existing trait system.

\textbf{Estimated scope:} \textasciitilde{}3 library files, \textasciitilde{}30 tests.

\begin{verbatim}
;; in lib/prologos/core/complex.prologos
type Complex := (complex Real Real)

spec complex-add : Complex -> Complex -> Complex
spec complex-mul : Complex -> Complex -> Complex
spec complex-conjugate : Complex -> Complex
spec complex-magnitude-squared : Complex -> Real

impl Semiring Complex
  zero := (complex 0.0 0.0)
  one  := (complex 1.0 0.0)
  add  := complex-add
  mul  := complex-mul
\end{verbatim}
\subsection{7.2 Phase T2: AmplitudeCell}
\label{sec:org7446fc6}

\textbf{Deliverables:}
\begin{itemize}
\item \texttt{AmplitudeCell} type: Map V Complex
\item Lattice instance with pointwise complex addition merge
\item Basic interference tests
\end{itemize}

\textbf{Dependencies:} Phase T1.

\textbf{Estimated scope:} \textasciitilde{}2 library files, \textasciitilde{}25 tests.

\textbf{Key tests:}
\begin{itemize}
\item Two-path interference: amplitudes a, b → probability |a+b|²
\item Destructive interference: a = -b → probability 0
\item Constructive interference: a = b → probability 4|a|²
\item Merge monotonicity under information ordering
\end{itemize}
\subsection{7.3 Phase T3: TensorCell}
\label{sec:orgfa836d7}

\textbf{Deliverables:}
\begin{itemize}
\item \texttt{TensorCell} type: Map (Tuple Vs) Complex
\item Tensor product construction for combining cells
\item Partial trace (marginalization) operation
\item Schmidt decomposition for entanglement analysis
\end{itemize}

\textbf{Dependencies:} Phase T2.

\textbf{Estimated scope:} \textasciitilde{}3 library files, \textasciitilde{}40 tests.

\textbf{Key tests:}
\begin{itemize}
\item Product state: TensorCell factors into independent AmplitudeCells
\item Bell state: TensorCell does NOT factor (entanglement)
\item Partial trace produces correct marginals
\item Schmidt rank correctly identifies entanglement degree
\end{itemize}
\subsection{7.4 Phase T4: Unitary Propagators}
\label{sec:orga794a54}

\textbf{Deliverables:}
\begin{itemize}
\item Matrix type with complex entries
\item Unitary propagator type in \texttt{propagator.rkt}
\item Composition and tensor product of unitary operators
\item Dagger (adjoint) operation
\item Standard gates: Hadamard, Phase, CNOT, Pauli X/Y/Z
\end{itemize}

\textbf{Dependencies:} Phase T3.

\textbf{Estimated scope:} \textasciitilde{}4 files (1 library, 1 propagator extension, 2 test), \textasciitilde{}50 tests.

\textbf{Key implementation:} Extend \texttt{propagator.rkt} with a new propagator variant that
applies matrix multiplication rather than lattice join. The unitary propagator is
\emph{not} added to the standard worklist; instead, it is fired during the unitary
stratum (Layer 3) of the four-layer evaluation cycle.

\begin{verbatim}
;; Extension to propagator.rkt (pseudo-code)
(struct unitary-propagator
  (matrix        ;; Matrix Complex n n (the unitary operator)
   input-cell    ;; cell-id (TensorCell to transform)
   output-cell   ;; cell-id (TensorCell to write result)
   stratum)      ;; Nat (which unitary stratum this belongs to)
  #:transparent)

(define (fire-unitary-propagator net uprop)
  (define input-val (net-cell-read net (unitary-propagator-input-cell uprop)))
  (define output-val (matrix-mul (unitary-propagator-matrix uprop) input-val))
  (net-cell-write net (unitary-propagator-output-cell uprop) output-val))
\end{verbatim}
\subsection{7.5 Phase T5: Measurement Stratum}
\label{sec:org5bfccdc}

\textbf{Deliverables:}
\begin{itemize}
\item Projection propagators (project TensorCell onto subspace)
\item Born rule implementation (amplitudes → probabilities)
\item ATMS integration for measurement outcomes (each outcome = assumption)
\item Four-layer evaluation cycle: monotonic → ATMS → unitary → measurement
\end{itemize}

\textbf{Dependencies:} Phase T4, existing ATMS.

\textbf{Estimated scope:} \textasciitilde{}3 files, \textasciitilde{}40 tests.

\textbf{Key implementation:} The four-layer cycle is orchestrated by a new
\texttt{tensor-aware-solve} function that extends \texttt{stratified-solve-goal}:

\begin{verbatim}
;; Orchestration (pseudo-code)
(define (tensor-aware-solve net goals)
  ;; Layer 1: Run monotonic propagation to quiescence
  (define net1 (run-to-quiescence net))

  ;; Layer 2: ATMS choice points (if any)
  (define net2 (atms-select-worldview net1))

  ;; Layer 3: Apply all pending unitary propagators
  (define net3 (fire-all-unitary-propagators net2))

  ;; Layer 4: Apply all pending measurements
  ;; Each measurement creates ATMS assumptions for outcomes
  (define net4 (fire-all-measurement-propagators net3))

  ;; If new information was produced, repeat
  (if (net-changed? net4 net)
      (tensor-aware-solve net4 goals)
      net4))
\end{verbatim}
\subsection{7.6 Phase T6: DensityCell and CPM}
\label{sec:org110f8e2}

\textbf{Deliverables:}
\begin{itemize}
\item \texttt{DensityCell} type: Map (Pair V V) Complex
\item Completely positive map propagators
\item Decoherence modeling (GKSL master equation discretization)
\item Classical-quantum hybrid networks
\end{itemize}

\textbf{Dependencies:} Phase T5.

\textbf{Estimated scope:} \textasciitilde{}4 files, \textasciitilde{}35 tests.

This phase enables the open quantum systems approach used in cognitive modeling.
The GKSL master equation dρ/dt = -i[H,ρ] + Σ\textsubscript{k}(L\textsubscript{k} ρ L\textsubscript{k}† - ½\{L\textsubscript{k}†L\textsubscript{k}, ρ\})
is discretized as a completely positive map applied to a DensityCell.
\subsection{7.7 Phase T7: ZX Rewrite Engine (Research Phase)}
\label{sec:org474c026}

\textbf{Deliverables:}
\begin{itemize}
\item ZX diagram representation as a data structure
\item Core rewrite rules (spider fusion, bialgebra, copy, color change)
\item Diagram-to-propagator-network extraction
\item Automated network simplification via rewriting
\end{itemize}

\textbf{Dependencies:} Phase T4.

\textbf{Estimated scope:} Research-level; \textasciitilde{}6 files, open-ended tests.

This is a research phase --- the deliverable is a prototype, not production
infrastructure. The goal is to determine whether ZX-calculus optimization
provides meaningful performance improvements for tensor-aware propagator
networks.
\subsection{7.8 Phase T8: Surface Syntax}
\label{sec:org5b4c114}

\textbf{Deliverables:}
\begin{itemize}
\item \texttt{solve-with :semiring complex} for quantum-like queries
\item \texttt{\textasciicircum{}} weight annotation extended to complex values
\item Unitary operator definition syntax
\item Measurement syntax
\item Entanglement construction syntax
\end{itemize}

\textbf{Dependencies:} Phases T5-T6.

\textbf{Estimated scope:} Parser extensions, elaborator extensions, \textasciitilde{}50 tests.

\begin{verbatim}
;; Quantum-like decision model in Prologos surface syntax
defr belief-state [?agent:Agent ?state:(AmplitudeCell Choice)]
  || "alice" (amplitude-cell
               (pair "cooperate" (complex 0.6 0.3))
               (pair "defect"    (complex 0.4 -0.2))) ^1.0

;; Unitary evolution of belief under social influence
def social-influence : Matrix Complex 2 2 :=
  (unitary-from-angle 0.3)  ;; rotation by 0.3 radians

;; Query with measurement
let result := (solve-with :semiring complex
  (evolve belief-state "alice" social-influence)
  (measure belief-state "alice" ?outcome))
\end{verbatim}
\section{8. Connection to Quantum Social and Economic Models}
\label{sec:org4938a41}

A companion document ("Quantum-Like Models for Social, Economic, and Cognitive
Systems") provides comprehensive background on the research landscape. Here we
summarize how the TensorCell architecture enables concrete modeling of these
phenomena.
\subsection{8.1 Decision Modeling}
\label{sec:org8758d31}

Yukalov-Sornette Quantum Decision Theory decomposes behavioral probability into
a utility factor (classical) and an attraction factor (quantum interference).
In the propagator network:

\begin{itemize}
\item Each prospect is a basis element in an AmplitudeCell.
\item The decision-maker's state is a superposition of prospects.
\item The attraction factor emerges as the interference term when amplitudes are
squared to produce choice probabilities.
\item The Quarter Law (average |q| = 1/4) provides a testable prediction.
\end{itemize}

\begin{verbatim}
;; Modeling a decision under uncertainty
defr prospect-state [?agent:Agent ?state:(AmplitudeCell Prospect)]

;; Interference produces non-classical choice probabilities
;; p(A) = |amplitude(A)|² + interference_term(A)
;; where interference_term arises from off-diagonal coherences
\end{verbatim}
\subsection{8.2 Game-Theoretic Modeling}
\label{sec:orgf97fdd6}

The Eisert-Wilkens-Lewenstein (EWL) protocol maps directly to TensorCells:

\begin{itemize}
\item Two players' strategy spaces form a tensor product: TensorCell [Strategy Strategy].
\item The entangling gate J(γ) is a unitary propagator applied to the joint state.
\item Players' individual strategies U\textsubscript{A}, U\textsubscript{B} are unitary propagators applied to
their respective subsystems.
\item Payoffs are computed by measurement of the final state.
\item Nash equilibrium search is propagation to quiescence under the payoff constraints.
\end{itemize}

The quantum Prisoner's Dilemma --- where entangled strategies produce a Pareto-
optimal Nash equilibrium that classical strategies cannot reach --- is expressible
as a propagator network with TensorCells.
\subsection{8.3 Financial Modeling}
\label{sec:org10d55b5}

Haven's connection between the Black-Scholes equation and the Schrödinger equation
suggests modeling financial instruments as amplitude states:

\begin{itemize}
\item Asset states in superposition until "measured" (traded).
\item The arbitrage parameter ℏ measures market efficiency.
\item Interference between bullish and bearish sentiment produces non-classical
price dynamics.
\item The tropical semiring still works for optimization; the complex semiring
adds interference effects.
\end{itemize}
\subsection{8.4 Social Network Modeling}
\label{sec:org7e33fb0}

Opinion dynamics via quantum walks:

\begin{itemize}
\item Each agent's belief is an AmplitudeCell over opinion options.
\item Social influence is a unitary rotation toward a neighbor's state.
\item Measurement = when an agent must act (vote, purchase, commit).
\item Network-level entanglement models correlated belief updates between connected
agents.
\item The Lindblad-Kossakowski equation (implemented via DensityCell) models
decoherence as consensus emergence.
\end{itemize}
\subsection{8.5 Survey and Polling Models}
\label{sec:org34688df}

The Wang-Busemeyer QQ equality --- a parameter-free prediction confirmed on
70+ national surveys --- is naturally expressed as a constraint on an
AmplitudeCell:

\begin{verbatim}
;; The QQ equality as a propagator constraint
;; p(A,B) + p(¬A,¬B) - p(B,A) - p(¬B,¬A) = 0
;; This is automatically satisfied by any quantum probability model
;; Violation would indicate the model is not quantum-like
defr qq-equality [?agent:Agent ?questionA ?questionB]
  :constraint
  &> (let pAB := (sequential-probability agent questionA questionB))
     (let pAnBn := (sequential-probability agent (not questionA) (not questionB)))
     (let pBA := (sequential-probability agent questionB questionA))
     (let pBnAn := (sequential-probability agent (not questionB) (not questionA)))
     (assert-equal (add pAB pAnBn) (add pBA pBnAn))
\end{verbatim}
\section{9. Research Questions and Open Problems}
\label{sec:orga81269b}

\subsection{9.1 Convergence Guarantees for the Four-Layer Cycle}
\label{sec:org8b3b433}

The monotonic layer converges by Tarski's fixed-point theorem (every monotone
function on a complete lattice has a least fixed point). The stratified layer
converges by the stratification guarantee (each stratum is finite). But the
four-layer cycle (monotonic → ATMS → unitary → measurement → monotonic → \ldots{})
may not converge if measurements produce information that triggers new unitary
evolutions that produce new measurements.

\textbf{Question:} Under what conditions does the four-layer cycle converge? Is there
a categorical characterization of convergent tensor-aware propagator networks?
\subsection{9.2 Scheduling and Measurement Order}
\label{sec:org787766b}

In quantum mechanics, the order of measurements matters (non-commutative
projections). In the propagator network, the order in which measurement
propagators fire may affect the result.

\textbf{Question:} How should the propagator scheduler handle measurement ordering?
Can the ATMS be extended to explore all measurement orderings, or is there a
canonical ordering derivable from the network topology?
\subsection{9.3 Compositionality}
\label{sec:org5659c5e}

How do tensor-aware sub-networks compose? If two sub-networks share a TensorCell,
their joint behavior should be well-defined. The compact closed structure provides
composition via the categorical trace, but the interaction with ATMS worldviews
and stratification needs analysis.

\textbf{Question:} Is there a compositionality theorem for tensor-aware propagator
networks? Under what conditions can a network be decomposed into independently
evaluable sub-networks?
\subsection{9.4 Tensor Network Contraction for Performance}
\label{sec:org9048f3b}

Tensor product spaces grow exponentially with the number of variables (2\textsuperscript{n} for
n binary variables). Tensor network contraction methods (Biamonte et al., 2018)
provide efficient algorithms for evaluating certain tensor network structures
without materializing the full tensor.

\textbf{Question:} Can the propagator network topology guide tensor network contraction
strategies? Can the widening operator for TensorCells (truncating small
amplitudes) be made systematic via the Widenable trait?
\subsection{9.5 Type-Level Entanglement Tracking}
\label{sec:orge00995f}

Prologos's dependent type system could track entanglement structure at the type
level. A \texttt{TensorCell \{[A B]\}} that is guaranteed separable could have a
different type than one that may be entangled, enabling compile-time verification
of entanglement requirements.

\textbf{Question:} What is the appropriate type-level representation of entanglement?
Can Schmidt rank be tracked in the type system? Can the fibered category
structure (Beyond Prolog §7.3) accommodate entanglement as a fiber property?
\subsection{9.6 Connection to Quantum Error Correction}
\label{sec:orgf254c5f}

Quantum error correction codes protect information against decoherence by
encoding logical qubits into physical qubits. The stabilizer formalism
(Gottesman) has a clean algebraic structure that may connect to the propagator
network's constraint propagation.

\textbf{Question:} Can quantum error correction be expressed as propagator constraints
on TensorCells? Would the ZX-calculus rewrite engine (Phase T7) enable
automated discovery of error-correcting codes?
\subsection{9.7 The ZX-Calculus Optimization Question}
\label{sec:org5d8f3eb}

The ZX-calculus provides a complete rewrite system for qubit quantum mechanics.
Whether this translates to practical optimization of tensor-aware propagator
networks is an open question.

\textbf{Question:} Does ZX-calculus optimization provide meaningful speedups for the
specific class of tensor-aware propagator networks that arise from quantum-like
social/economic models? Or are these models structurally simpler than general
quantum circuits, rendering ZX optimization unnecessary?
\subsection{9.8 Non-Markovian Extensions}
\label{sec:org99090b4}

The GKSL master equation assumes Markovian dynamics (no memory effects). Human
cognition and social systems have strong memory effects. Non-Markovian quantum
dynamics (Nakajima-Zwanzig, time-convolutionless approaches) are more appropriate
but more complex.

\textbf{Question:} Can non-Markovian dynamics be expressed within the propagator network
framework? The persistent (immutable) nature of Prologos's propagator network
means historical states are naturally available --- can this be exploited for
memory effects?
\section{10. References}
\label{sec:orgb7aa55c}

\subsection{Categorical Foundations}
\label{sec:org5e9171d}
\begin{itemize}
\item Abramsky, S. \& Coecke, B. (2004). A categorical semantics of quantum protocols. \emph{LICS'04}.
\item Coecke, B. \& Duncan, R. (2008/2011). Interacting quantum observables. \emph{New J. Phys.} 13, 043016.
\item Coecke, B. \& Kissinger, A. (2017). \emph{Picturing Quantum Processes}. Cambridge UP.
\item Hadzihasanovic, A., Ng, K.F. \& Wang, Q. (2018). Two complete axiomatisations of pure-state qubit quantum computing. \emph{LMCS}.
\item Kelly, G.M. (1982). \emph{Basic Concepts of Enriched Category Theory}. Cambridge UP.
\item Lawvere, F.W. (1973). Metric spaces, generalized logic, and closed categories. \emph{Rend. Sem. Mat. Fis. Milano} 43.
\item Mulvey, C.J. (1986). \&. \emph{Rend. Circ. Mat. Palermo} Suppl. 12, 99-104.
\item Selinger, P. (2007). Dagger compact closed categories and completely positive maps. \emph{ENTCS} 170, 139-163.
\item van de Wetering, J. (2020). ZX-calculus for the working quantum computer scientist. \emph{arXiv:2012.13966}.
\end{itemize}
\subsection{Quantales and Enriched Categories}
\label{sec:orgf46876a}
\begin{itemize}
\item Hofmann, D., Seal, G.J. \& Tholen, W. (2014). \emph{Monoidal Topology}. Cambridge UP.
\item Kruml, D., Pelletier, J.W. \& Resende, P. (2003). On quantales and spectra of C*-algebras. \emph{Appl. Cat. Struct.} 11, 543-560.
\item Rosenthal, K.I. (1990). \emph{Quantales and their Applications}. Pitman.
\item Stubbe, I. (2013). An introduction to quantaloid-enriched categories. \emph{Fuzzy Sets and Systems}.
\end{itemize}
\subsection{Propagator Networks}
\label{sec:org4a3177a}
\begin{itemize}
\item Radul, A. (2009). Propagation networks: A flexible and expressive substrate for computation. \emph{PhD thesis, MIT}.
\item Radul, A. \& Sussman, G.J. (2009). The art of the propagator. \emph{MIT CSAIL TR-2009-002}.
\item Arntzenius, M. Semilattices and their tensor products. \emph{Blog post}.
\end{itemize}
\subsection{Stratified and Distributed Computation}
\label{sec:orgadbf96a}
\begin{itemize}
\item Alvaro, P. \& Hellerstein, J.M. (2019). Keeping CALM: When distributed consistency is easy. \emph{arXiv:1901.01930}.
\item Biamonte, J. et al. (2018). Tensor network contraction and the counting of CSPs. \emph{arXiv:1805.00475}.
\item Conway, N. et al. (2012). Logic and lattices for distributed programming. \emph{SoCC'12}.
\item Kuper, L. et al. (2013). LVars: Lattice-based data structures for deterministic parallelism. \emph{FHPC'13}.
\end{itemize}
\subsection{Quantum Social Science (see companion document for full bibliography)}
\label{sec:org3948b53}
\begin{itemize}
\item Busemeyer, J.R. \& Bruza, P.D. (2012). \emph{Quantum Models of Cognition and Decision}. Cambridge UP.
\item Eisert, J., Wilkens, M. \& Lewenstein, M. (1999). Quantum games and quantum strategies. \emph{Phys. Rev. Lett.} 83, 3077.
\item Haven, E. \& Khrennikov, A. (2013). \emph{Quantum Social Science}. Cambridge UP.
\item Wang, Z. \& Busemeyer, J.R. (2013). A quantum question order model supported by empirical tests. \emph{Topics in Cognitive Science} 5, 689-710.
\item Yukalov, V.I. \& Sornette, D. (2011). Decision theory with prospect interference and entanglement. \emph{Theory and Decision} 70, 283-328.
\end{itemize}
\subsection{Prologos Internal Documents}
\label{sec:org2f25f8b}
\begin{itemize}
\item Beyond Prolog: Frontiers for the Prologos Relational Language (2026-03-16).
\item Logic Engine Design: Comprehensive Implementation Guide (2026-02-24).
\item Towards a General Logic Engine on Propagators (2026-02-24).
\item Galois Connections \& Abstract Interpretation on the Logic Engine (2026-02-27).
\item Stratified Evaluation for the Logic Engine (2026-02-26).
\item Quantum-Like Models for Social, Economic, and Cognitive Systems (companion, 2026-03-16).
\end{itemize}
\end{document}
